\section{\texorpdfstring{第~\ref{chap:riemannian-manifolds}~章}{Riemann 流形}}\label{sec:solutions-riemannian-manifolds}
\subsection{\texorpdfstring{第~\ref{sec:riemannian-metrics}~节}{Riemann 度量}}\label{subsec:solutions-riemannian-metrics}

\begin{description}[font=\normalfont,leftmargin=2em,labelwidth=1.5em]
\RGsolution{ex:quotient-riemannian-structures}{quotient-riemannian-structures}
\begin{enumerate}[label=(\alph*),ref=\alph*,leftmargin=*]
\item 在这种情形下, 自然投影 $\pi:M\to M/G$ 是覆盖映射 (因而是局部微分同胚). 对每个 $r\in M/G$, 可以选取 $p\in\pi^{-1}(r)$ 及其邻域 $U\ni p$, 使得 $\pi|_U$ 是到其像的微分同胚. 然后在 $\pi(U)$ 上定义度量 $h:=((\pi|_U)^{-1})^*g$. 为证明这个定义不依赖于 $p$ 的选取, 注意 $\pi^{-1}(r)$ 中任意另一个点 $q$ 都具有形式 $q=f(p)$, 其中 $f:M\to M$ 是 $G$ 中某个元素的作用所对应的等距映射. 此时 $V:=f(U)$ 是 $q$ 的一个邻域, $\pi|_V$ 是微分同胚, 且 $\pi|_U=\pi|_V\circ f$. 因此,
\begin{align*}
  ((\pi|_U)^{-1})^*g
  &=(f^{-1}\circ(\pi|_V)^{-1})^*g
  =((\pi|_V)^{-1})^*(f^{-1})^*g\\
  &=((\pi|_V)^{-1})^*g,
\end{align*}
因为 $f$ (因而 $f^{-1}$) 是等距映射.
\item 在 $(M,g)$ 为带 Euclid 度量的 $\mathbb{R}^n$, 且 $\mathbb{Z}^n$ 以平移 (这是 Euclid 度量的等距映射) 作用于 $\mathbb{R}^n$ 的情形下, 采用上述方法, 就可以定义 $T^n=\mathbb{R}^n/\mathbb{Z}^n$ 上的平坦方形度量. 特别地, $T^n$ 的局部几何与 $\mathbb{R}^n$ 的局部几何无从区分. 有 $\operatorname{vol}(T^n)=1$, 因为若令 $U:=(0,1)^n$, 则 $\pi|_U:U\to T^n$ 是一个微分同胚, 它覆盖 $T^n$ 中除一个零测集之外的部分. 注意, $\mathbb{R}^n$ 的每一组基都确定 $\mathbb{Z}^n$ 的一个作用, 即按基向量的整数倍作平移, 所得商空间与 $T^n$ 微分同胚. 由这些作用在 $T^n$ 上得到的度量一般不同 (特别是, 相应的 $T^n$ 的体积不一定为 1).
\item 在 $(M,g)$ 为带标准度量的 $S^n$, 且 $\mathbb{Z}_2$ 通过对径映射 (这是标准度量的等距映射) 作用于 $S^n$ 的情形下, 采用上述方法, 就可以定义 $\mathbb{R}P^n=S^n/\mathbb{Z}_2$ 上的标准度量. 特别地, $\mathbb{R}P^n$ 的局部几何与 $S^n$ 的局部几何无从区分. 还应注意, $\operatorname{vol}(\mathbb{R}P^n)=\frac12\operatorname{vol}(S^n)$, 因为若 $U$ 是一个半球, 则 $\pi|_U:U\to\mathbb{R}P^n$ 是一个微分同胚, 它覆盖 $\mathbb{R}P^n$ 中除一个零测集之外的部分.
\end{enumerate}

\RGsolution{ex:left-invariant-metric}{left-invariant-metric}
\begin{enumerate}[label=(\alph*),ref=\alph*,leftmargin=*]
\item 若 $g$ 左不变, 则对所有 $x\in G$ 和所有 $v,w\in T_xG$, 必有
\[
  \langle v,w\rangle_x
  =\langle(\mathrm{d}L_{x^{-1}})_xv,(\mathrm{d}L_{x^{-1}})_xw\rangle_e.
\]
所以只需证明此式确实定义了 $G$ 上的一个左不变度量. 容易验证, 映射
\[
  G\times G\ni(x,y)\mapsto x^{-1}y=L_{x^{-1}}y\in G
\]
的光滑性蕴含映射
\[
  G\times TG\ni(x,v)\mapsto(\mathrm{d}L_{x^{-1}})_xv\in TG
\]
的光滑性. 因此上式定义了 $G$ 上的光滑张量场 $g$. $g$ 对称且正定也是显然的. 只剩下证明 $g$ 左不变, 即对所有 $v,w\in T_xG$ 及所有 $x,y\in G$,
\[
  \langle(\mathrm{d}L_y)_xv,(\mathrm{d}L_y)_xw\rangle_{yx}=\langle v,w\rangle_x.
\]
实际上,
\begin{align*}
  &\langle(\mathrm{d}L_y)_xv,(\mathrm{d}L_y)_xw\rangle_{yx}\\
  &\quad=\left\langle(\mathrm{d}L_{(yx)^{-1}})_{yx}(\mathrm{d}L_y)_xv,
                     (\mathrm{d}L_{(yx)^{-1}})_{yx}(\mathrm{d}L_y)_xw\right\rangle_e\\
  &\quad=\left\langle\bigl(\mathrm{d}(L_{x^{-1}y^{-1}}\circ L_y)\bigr)_xv,
                     \bigl(\mathrm{d}(L_{x^{-1}y^{-1}}\circ L_y)\bigr)_xw\right\rangle_e\\
  &\quad=\langle(\mathrm{d}L_{x^{-1}})_xv,(\mathrm{d}L_{x^{-1}})_xw\rangle_e=\langle v,w\rangle_x.
\end{align*}
所以 Lie 代数 $\mathfrak g=T_eG$ 上的任意内积都确定 $G$ 上的一个左不变度量.
\item 回顾, 每个四元数 $q\in\mathbb{R}SU(2)$ 都可写成
\[
  q=a1+bi+cj+dk,
\]
其中 $a,b,c,d\in\mathbb{R}$, 而 $1,i,j,k\in SU(2)$ 在第~\ref{chap:differentiable-manifolds}~章的 Exercise~\ref{ex:lie-groups}(\ref{ex:quaternions-and-covering-of-rotation-group}) 中给出; $SU(2)$ 是 Euclid 长度为 1 的四元数所成的集合. 由恒等式 $i^2=j^2=k^2=ijk=-1$, 容易验证
\[
  qq^*=q^*q=(a^2+b^2+c^2+d^2)1.
\]
此外, 用这组基将 $\mathbb{R}SU(2)$ 与 $\mathbb{R}^4$ 等同, 则 Euclid 内积可写成
\[
  \langle q_1,q_2\rangle=\operatorname{Re}(q_1q_2^*)=\operatorname{Re}(q_1^*q_2),
\]
其中 $\operatorname{Re}(q)=a$. 若 $v\in SU(2)$ 是单位四元数, 即 $v^*v=1$, 则
\[
  \langle vq_1,vq_2\rangle=\operatorname{Re}(q_1^*v^*vq_2)
  =\operatorname{Re}(q_1^*q_2)=\langle q_1,q_2\rangle.
\]
所以乘以单位四元数保持 Euclid 内积. 将其限制在 $SU(2)$ 的切向量上, 可知 $SU(2)$ 上的标准度量左不变.
\item $\mathcal{M}_{n\times n}\cong\mathbb{R}^{n^2}$ 上的 Euclid 内积由
\[
  \langle A,B\rangle=\operatorname{tr}(AB^t)
\]
给出. 因此, 若 $S\in O(n)$, 则
\[
  \langle SA,SB\rangle=\operatorname{tr}(SAB^tS^t)
  =\operatorname{tr}(S^tSAB^t)=\operatorname{tr}(AB^t)=\langle A,B\rangle.
\]
将其限制在 $O(n)$ 的切向量上, 可知由 $\mathcal{M}_{n\times n}\cong\mathbb{R}^{n^2}$ 上的 Euclid 度量在 $O(n)$ 上诱导的度量左不变.
\end{enumerate}
\end{description}
